%% 01_introduction.tex
%% Section 1: Introduction

\section{Introduction}
\label{sec:introduction}

Indonesia occupies a singular position in the global climate-finance landscape. As the world's largest thermal coal exporter, shipping over 500 million tonnes annually and supplying roughly one-quarter of global seaborne trade, its coal sector sits at the epicentre of the unburnable carbon problem \citep{McGlade2015, Welsby2021}. Yet this exposure extends far beyond the mining sector itself. Indonesian banks hold substantial credit portfolios tied to coal mining and related activities, while the sovereign budget depends materially on coal royalties and corporate tax revenues. This triple entanglement, what we term the \textit{coal-banking-sovereign nexus}, means that a disorderly climate transition could propagate from stranded coal assets through the financial system to the broader macroeconomy in ways that are mutually reinforcing and potentially destabilising.

The concept of stranded assets, broadly defined as fossil fuel reserves and associated infrastructure that lose economic value under climate-constrained pathways, has attracted growing scholarly attention since the seminal contributions of \citet{McGlade2015} and the Carbon Tracker Initiative \citep{CarbonTracker2022}. \citet{Welsby2021} estimate that 89\% of global coal reserves must remain unextracted to meet the 1.5\textdegree C target, with Southeast Asian coal particularly vulnerable due to its relatively high extraction costs and distance from end markets. The financial implications of such stranding are profound: \citet{Mercure2018} estimate global wealth losses of \$1--4 trillion from stranded fossil fuel assets, while \citet{Dietz2016} place the climate value-at-risk of global financial assets at \$2.5 trillion under business-as-usual warming.

Despite this growing body of work, a critical gap remains. Existing studies of stranded assets have largely adopted either a \textit{top-down} perspective (estimating global or regional reserves that must stay underground) or a \textit{sector-specific} focus on coal, oil, or gas without tracing the financial transmission channels. Studies of climate stress testing in the banking sector, such as those by \citet{Battiston2017}, \citet{Vermeulen2018}, and the ECB \citep{ECB2022}, have advanced our understanding of how climate risk enters bank balance sheets, but these analyses are predominantly focused on European financial systems. The macro-financial transmission of transition risk, examined theoretically by \citet{Carattini2023} and \citet{Monasterolo2020}, has not been empirically grounded in firm-level stranding data. For Indonesia specifically, the Institute for Essential Services Reform \citep{IESR2023} provides the most detailed country-level assessment to date, but its banking analysis remains qualitative and does not quantify macroeconomic spillovers.

This paper fills these gaps by developing the first integrated, bottom-up framework linking coal stranded asset risk to banking sector vulnerability in Indonesia. Our approach is distinctive in two respects. First, we construct firm-level stranded asset valuations for 35 listed coal companies on the Indonesia Stock Exchange (IDX), using a supply cost curve methodology combined with scenario-contingent discounted cash flow analysis and Monte Carlo uncertainty quantification. Second, we build a bipartite bank-coal exposure network from granular loan data covering 38 listed banks and conduct a climate stress test that estimates credit losses and capital adequacy impacts at the individual bank level.

Our analysis yields several striking findings. Under a 2\textdegree C-aligned transition scenario, we estimate aggregate stranded values of \$79.7 billion (90\% confidence interval: \$72.7--88.0 billion), rising to \$128.3 billion (90\% CI: \$117.1--141.0 billion) under the 1.5\textdegree C pathway. These losses propagate through the banking system, generating total bank credit losses of \$1.8 billion under the 2\textdegree C scenario and \$3.1 billion under the 1.5\textdegree C scenario. The stress test reveals that while no bank breaches the OJK minimum capital requirement, the most exposed D-SIB sees its capital adequacy ratio decline by up to 1.00 percentage point, with losses absorbing 42\% of annual earnings. The bipartite network structure further reveals that the high degree of overlap between lenders, combined with correlated coal default risk, limits the diversification available to banks within the sector, underscoring the systemic dimension of concentrated coal lending.

The relevance of this analysis has been sharpened by Indonesia's entry into the Just Energy Transition Partnership (JETP) in November 2022. The JETP mobilises \$20 billion in public and private finance to accelerate Indonesia's energy transition, centred on capping power sector emissions at 290~Mt CO\textsubscript{2} by 2030 and achieving net-zero electricity by 2050 \citep{JETP2022}. Delivering on these commitments requires retiring a substantial share of the existing coal-fired power fleet and redirecting private capital toward renewables. Yet both objectives place the banking sector at the centre of the transition: banks must absorb losses on stranded coal assets while simultaneously scaling up green finance. Our analysis demonstrates that a rigorous assessment of banking sector vulnerabilities is a necessary precondition for credible JETP implementation, not an afterthought. Without such an assessment, JETP's financing plan rests on an incomplete picture of the financial system's resilience.

Beyond Indonesia, there is a growing international demand for country-level climate stress tests calibrated to emerging market conditions. The Network for Greening the Financial System (NGFS), whose membership now encompasses over 130 central banks and supervisors, has called explicitly for standardised scenario analysis and climate stress testing at the country level, with particular urgency for economies whose financial systems are heavily exposed to fossil fuel transition risk \citep{NGFS2022}. The Task Force on Climate-related Financial Disclosures (TCFD) similarly emphasises scenario analysis as a core disclosure requirement, creating a regulatory expectation that financial institutions must quantify, not merely describe, their climate risk exposures. Indonesia's OJK and Bank Indonesia have taken initial steps in this direction through the Sustainable Finance Roadmap and TCFD-aligned disclosure circulars, but a firm-level, bottom-up empirical foundation for these exercises has been absent. This paper provides that foundation for the coal-banking nexus specifically.

The \textit{coal-banking-sovereign nexus} framing that organises our analysis reflects a structural feature of Indonesia's political economy that distinguishes it from most countries studied in the existing climate finance literature. The nexus operates through three mutually reinforcing channels: coal royalties and taxes underpin sovereign fiscal capacity; state banks and private lenders hold material coal credit portfolios; and the coal sector's dominant position in the export base means that transition-induced revenue losses reverberate through the broader economy. Indonesia is arguably the most important case study for this nexus globally: no other country combines such a large share of global seaborne coal trade, such a bank-centred financial system, and such an explicit commitment to the JETP transition framework. Understanding the nexus in the Indonesian context thus provides both a critical country-level assessment and a methodological template transferable to other coal-dependent emerging economies, including South Africa, India, Colombia, and the Philippines.

These results carry important implications for climate-financial policy in commodity-dependent emerging economies. Our framework provides regulators, including the Financial Services Authority (OJK) and Bank Indonesia (BI), with a quantitative toolkit for assessing the systemic dimensions of transition risk. For OJK specifically, the results call for immediate incorporation of climate transition risk into the supervisory stress testing programme, implementation of climate-adjusted capital requirements for coal-exposed banks, and enhanced borrower-level fossil fuel disclosure standards aligned with NGFS guidance. For Bank Indonesia, the findings underscore the case for developing green refinancing facilities that incentivise banks to redirect lending from coal to renewable energy, echoing the green monetary policy instruments adopted by the ECB and the People's Bank of China. More broadly, the coal-banking-sovereign nexus we document is likely present in other coal-dependent economies such as South Africa, India, Colombia, and the Philippines, making our methodology and findings of wider relevance for the international climate finance architecture.

The remainder of this paper is organised as follows. Section~\ref{sec:literature} reviews the relevant literatures on stranded assets, climate stress testing, and the Indonesia-specific context. Section~\ref{sec:background} provides institutional background on Indonesia's coal sector, banking system, fiscal dependence on coal revenues, and the evolving regulatory landscape. Section~\ref{sec:data} describes our data sources and presents descriptive statistics. Sections~\ref{sec:model1} and~\ref{sec:model2} develop and present results from our two interconnected models. Section~\ref{sec:discussion} discusses policy implications, compares our estimates with prior work, and acknowledges limitations. Section~\ref{sec:conclusion} concludes.
